\section{A Unified Framework of LLM Training Algorithms}

\subsection{Pre-Training And Distillation Algorithms}

In pre-training and Supervised Fine-Tuning (SFT), minimizing the standard
cross-entropy loss is equivalent to minimizing the forward KL divergence
from the data distribution to the model distribution.
To make this precise, let $\mathcal{D}=\{x^{(i)}\}_{i=1}^N$ be a fixed,
nonempty dataset of token sequences, let
$p_{\mathcal{D}}(x)=N^{-1}\sum_{i=1}^N\mathbf{1}\{x^{(i)}=x\}$ be its
empirical distribution, and let $\pi_\theta(x)$ be the model probability
of sequence $x$ with parameter $\theta$. Then

\begin{align}
 D_{\mathrm{KL}}(p_{\mathcal{D}}\Vert\pi_\theta)
 &= \mathbb{E}_{x\sim p_{\mathcal{D}}}
    \left[\log\frac{p_{\mathcal{D}}(x)}{\pi_\theta(x)}\right] \notag = -\mathbb{E}_{x\sim p_{\mathcal{D}}}[\log\pi_\theta(x)]
    -H(p_{\mathcal{D}}) \notag\\
 &= \mathcal{L}_{\mathrm{CE}}(\theta)-H(p_{\mathcal{D}}),
 \label{eq:ce-forward-kl}
\end{align}

Here the empirical cross-entropy loss is
$\mathcal{L}_{\mathrm{CE}}(\theta)=-N^{-1}\sum_{i=1}^N
\log\pi_\theta(x^{(i)})$.
The data entropy $H(p_{\mathcal{D}})=-\mathbb{E}_{x\sim p_{\mathcal{D}}}
[\log p_{\mathcal{D}}(x)]$ is finite and independent of $\theta$ which has no effect on the gradient. This proves the equivalence.

For SFT, let $x$ denote a prompt and $y$ its target response, with fixed
empirical joint distribution $p_{\mathcal{D}}(x,y)$. The same decomposition
holds conditionally:

\begin{align}
  \mathcal{L}_{\mathrm{SFT}}(\theta) &= 
  -\mathbb{E}_{(x,y)\sim p_{\mathcal{D}}}
  [\log\pi_\theta(y\mid x)]\\
  &=\mathbb{E}_{x\sim p_{\mathcal{D}}}
 \left[D_{\mathrm{KL}}\bigl(p_{\mathcal{D}}(\cdot\mid x)
 \Vert\pi_\theta(\cdot\mid x)\bigr)\right]
 +H_{p_{\mathcal{D}}}(Y\mid X),
 \label{eq:sft-conditional-kl}
\end{align}

Under the autoregressive factorization,
$-\log\pi_\theta(y\mid c)=-\sum_{t=1}^{|y|}
\log\pi_\theta(y_t\mid c,y_{<t})$, this sequence-level objective is exactly the sum of token-level
cross-entropies with ground-truth prefixes.

In many applications the response could come from a teacher model $\pi_t$, instead of the dataset. Then the training objective of SFT would become minimizing the forward KL divergence between the teacher model and the student model. This yields the classical SFT-based Knowledge-Distillation (KD) algorithm~\citep{hinton_distilling_2015}, which is often simply called SFT as well:

\begin{align}
  \mathcal{J}_{\mathrm{KD}} &= -\mathbb{E}_{x\sim\mathcal{D}} [D_{\mathrm{KL}}\bigl(\pi_t (\cdot\mid x)\Vert \pi_\theta (\cdot\mid x)\bigr)] = \mathbb{E}_{x\sim\mathcal{D}, y\sim\pi_t (\cdot\mid x)} [\log \pi_\theta (y\mid x)] - H(\pi_t).
  \label{eq:kl-loss}
\end{align}

OPD recently emerged as a new algorithm for distillation tasks~\citep{thinking_machines_lab_-policy_2025}. It can exhibit superior performance than SFT in some cases, and received great attention from the industry~\citep{team_kimi_2026, deepseek-ai_deepseek-v4_2026, glm-5-team_glm-5_2026}. While in some cases this algorithm may also be unstable and failed to complete the distillation task~\citep{li_rethinking_2026}. The key difference between OPD and KD is that, OPD uses sequence-level reverse KL divergence as the training objective:

\begin{align}
  \mathcal{J}_{\mathrm{OPD}} &= -\mathbb{E}_{x\sim\mathcal{D}} [D_{\mathrm{KL}}\bigl(\pi_\theta (\cdot\mid x)\Vert \pi_t (\cdot\mid x)\bigr)]
  = \mathbb{E}_{x\sim\mathcal{D}, y\sim\pi_\theta (\cdot\mid x)} [\log \frac{\pi_t (y\mid x)}{\pi_\theta (y\mid x)}].
  \label{eq:opd-loss}
\end{align}

The on policy nature of OPD directly come from the requirement of reverse KL divergence. As the expectation of reverse KL is calculated under the student distribution, we should use the student to generate samples. We should note that the sequence-level KL objective directly yields sampled-token OPD. While there exist several variants of this OPD~\citep{li_rethinking_2026}, this sampled-token OPD is indeed adopted by the industry~\citep{team_kimi_2026} and share some important properties with other variants. the following analysis will primarily be built upon this algorithm.

\subsection{Reinforce Learning Algorithms}

All the algorithms methoned above has a clear definition of the target distribution. However, as another central component of post-training, Reinforcement Learning (RL) is not defined by such divergence. The objective of RL is to maximize the expected reward $r$ given to each possible action of the policy model, while for LLMs this action is typically a sequence generation. Its standard objective and gradient are~\citep{williams_simple_1992}:

\begin{equation}
  \mathcal{J}_{\mathrm{RL}} = \mathbb{E}_{x\sim\mathcal{D}, y\sim\pi_\theta (\cdot\mid x)} [r(x, y)],\qquad \nabla_\theta \mathcal{J}_{\mathrm{RL}} = \mathbb{E}_{x\sim\mathcal{D}, y\sim\pi_\theta (\cdot\mid x)} [r(x, y) \nabla_\theta \log \pi_\theta].
  \label{eq:std-rfalgo}
\end{equation}

% In this work we will show that RL algorithms can also be expressed in the form of distribution matching objective under 0-1 reward, and the analysis on the target distribution can further reveal some interesting properties of these algorithms.

In fact, this objective also implicitly introduces a target distribution. We first characterize an ideal policy update that improves expected reward
while remaining close to the current policy under a forward-KL restriction. (Similar result also holds for reverse-KL restriction)

\begin{theorem}[Reward reweighting under a forward-KL restriction]
\label{lem:reward-reweighting}
For a prompt $x$, the ideal update maximizing
$\mathbb E_{y\sim q}[r(x,y)]-\beta D_{\mathrm{KL}}(\pi_\theta\Vert q)$
over response distributions $q$, with KL penalty $\beta>0$, is
\begin{equation}
 q^*(y\mid x)
 =\frac{\beta\pi_\theta(y\mid x)}{\lambda_x-r(x,y)},
 \label{eq:reward-reciprocal-tilt}
\end{equation}
where $\lambda_x$ normalizes the distribution. Responses with the same
reward receive the same probability multiplier, preserving their
conditional distribution.
\end{theorem}
\noindent\hyperref[app:reward-reweighting]{Proof: Appendix~\ref*{app:reward-reweighting}.}

Under binary rewards, we can verify that this update preserves the success-conditioned distribution
$\pi_\theta^+(\cdot\mid x)
=\pi_\theta(\cdot\mid x,r(x,y)=1)$. If we use this distribution as the training target, matching
this conditional distribution with the forward KL gives the following identity.

\Needspace{9\baselineskip}
\begin{theorem}[Success-conditioned KL and the MaxRL limit]
\label{lem:conditional-kl-maxrl}
Let $\rho_\theta(x)=\mathbb E_{y\sim\pi_\theta(\cdot\mid x)}[r(x,y)]$
denote the success probability. Matching the success-conditioned target
$\pi_\theta^+$ by minimizing its forward KL to the current policy yields
\begin{equation}
 -\mathbb E_{x\sim\mathcal D}
 D_{\mathrm{KL}}\bigl(\pi_\theta^+(\cdot\mid x)
                  \Vert\pi_\theta(\cdot\mid x)\bigr)
 =\mathbb E_{x\sim\mathcal D}[\log\rho_\theta(x)]
 =\mathcal J_{\mathrm{MaxRL}}^{(\infty)}(\theta).
 \label{eq:conditional-kl-maxrl}
\end{equation}
\end{theorem}
\noindent\hyperref[app:conditional-kl-maxrl]{Proof: Appendix~\ref*{app:conditional-kl-maxrl}.}

This actually gives the limit objective function of MaxRL algorithm, which is previously derived from the maximum likelihood estimator of the predictive accuracy~\citep{tajwar_maximum_2026}. It is believed to have some superior performance on maintaining the diversity of the distribution, and the standard RL objective is the first-order approximation of MaxRL.

An auxiliary binary acceptance variable extends the conditional-KL identity to bounded nonbinary rewards on an augmented response space; Appendix~\ref{app:randomized-reward-extension} outlines this construction, while our main RL analysis remains focused on binary rewards.

\begin{table}[!htbp]
\centering
\caption[Training algorithms viewed as distribution matching]{Training algorithms viewed as distribution matching. Parameter-independent constants
are omitted. Conditional objectives are averaged over prompts
$x\sim\mathcal D$. \forwardKL{} places the target in the first argument;
\reverseKL{} places it in the second. RL and MaxRL use binary rewards and
$\rho_\theta(x)>0$, with the policy-dependent target $\pi_\theta^+$.
Algorithms are iterative if their target distribution depends on the current
model, and noniterative otherwise; \textcolor{frameworkpolicy}{$\pi_\theta$}
is highlighted in iterative targets.
$^{\dagger}$\,RL is its first-order approximation.}
\label{tab:distribution-matching}
\vspace{5pt}
\begingroup
\small
\setlength{\tabcolsep}{5pt}
\renewcommand{\arraystretch}{1.45}
\begin{tabularx}{\linewidth}{
 >{\raggedright\arraybackslash}p{0.13\linewidth}
 >{\centering\arraybackslash}X
 >{\centering\arraybackslash}p{0.235\linewidth}
 >{\raggedright\arraybackslash}p{0.16\linewidth}}
\specialrule{\heavyrulewidth}{0pt}{0pt}
\rowcolor{frameworktint}
\textcolor{frameworkink}{\textbf{Algorithm}} &
\textcolor{frameworkink}{\textbf{Objective function}} &
\textcolor{frameworkink}{\textbf{Target distribution}} &
\textcolor{frameworkink}{\textbf{Divergence}} \\
\algorithmgroup{Noniterative algorithms}
Pre-training &
$\mathbb E_{x\sim p_{\mathcal D}}[\log\pi_\theta(x)]$ &
$p_{\mathcal D}(x)$ & \forwardKL \\
SFT &
$\mathbb E_{y\sim p_{\mathcal D}(\cdot\mid x)}[\log\pi_\theta(y\mid x)]$ &
$p_{\mathcal D}(y\mid x)$ & \forwardKL \\
\addlinespace[3pt]
SFT-based KD &
$\mathbb E_{y\sim\pi_t(\cdot\mid x)}[\log\pi_\theta(y\mid x)]$ &
$\pi_t(y\mid x)$ & \forwardKL \\
OPD &
$\displaystyle\mathbb E_{y\sim\pi_\theta(\cdot\mid x)}
 \!\left[\log\frac{\pi_t(y\mid x)}{\pi_\theta(y\mid x)}\right]$ &
$\pi_t(y\mid x)$ & \reverseKL \\
\algorithmgroup{Iterative algorithms}
RL &
$\mathbb E_{y\sim\pi_\theta(\cdot\mid x)}[r(x,y)]$ &
$\displaystyle\targetpolicy^+=\frac{r(x,y)\targetpolicy(y\mid x)}{\rho_\theta(x)}$ &
\forwardKL$^{\dagger}$ \\
MaxRL &
$\log\mathbb E_{y\sim\pi_\theta(\cdot\mid x)}[r(x,y)]$ &
$\displaystyle\targetpolicy^+=\frac{r(x,y)\targetpolicy(y\mid x)}{\rho_\theta(x)}$ &
\forwardKL \\
\bottomrule
\end{tabularx}
\endgroup
\end{table}
\FloatBarrier

Table~\ref{tab:distribution-matching} summarizes these algorithms through
their objectives and target distributions.

% In the following sections we will primarily study the property of these algorithms under different conditions, with a focus on KD, OPD and RL with binary reward.
